\begin{thema}{Γ}
Στο σχήμα δίνονται η γραφική παράσταση $C_f$ μιας γνησίως αύξουσας συνάρτησης $f:[0, 5]\to [0,4]$ και η γραφική παράσταση $C_{f^{-1}}$ της αντίστροφης της.

\begin{center}
\begin{tikzpicture}[scale=0.75]
  \draw[help lines, gray!30] (-1,-1) grid (6,5);
  \draw[-latex] (-1.3,0) -- (6.3,0) node[right] {$x$};
  \draw[-latex] (0,-1.3) -- (0,5.3) node[above] {$y$};
  \foreach \x in {1,2,3,4,5}
    \draw (\x,0.1) -- (\x,-0.1) node[below,font=\footnotesize] {$\x$};
  \foreach \y in {1,2,3,4}
    \draw (0.1,\y) -- (-0.1,\y) node[left,font=\footnotesize] {$\y$};
  \draw[dashed, gray!50] (-1,-1) -- (5,5);
  % f: concave, below y=x
  \node[secondcolor!80!black] at (1,3.5) {$C_{f^{-1}}$};
  \filldraw[maincolor] (0,0) circle (2pt);
  \filldraw[maincolor] (1,1) circle (2pt) node[below right,font=\footnotesize] {$(1,1)$};
  \filldraw[maincolor] (3,2) circle (2pt) node[below right,font=\footnotesize] {$(3,2)$};
  \filldraw[maincolor] (5,4) circle (2pt) node[below right,font=\footnotesize] {$(5,4)$};
  % shaded between f and f^{-1}
  \fill[secondcolor!15, opacity=0.5, smooth, tension=0.4]
    plot coordinates {(1,1) (2,1.6) (3,2) (4,3)}
    -- plot[smooth, tension=0.4] coordinates {(4,3) (3,4) (2,3) (1.6,2) (1,1)} -- cycle;
  \draw[very thick, maincolor, smooth, tension=0.4]
    plot coordinates {(0,0) (1,1) (2,1.6) (3,2) (4,3) (5,4)};
  \node[maincolor] at (4.5,2) {$C_f$};
  % f^{-1}: convex, above y=x
  \draw[very thick, secondcolor!80!black, smooth, tension=0.4]
    plot coordinates {(0,0) (1,1) (1.6,2) (2,3) (3,4) (4,5)};
\end{tikzpicture}
\end{center}

\begin{erwthma}
  \item Να βρείτε τα $f^{-1}(2)$, $f^{-1}(4)$ και $(f \circ f^{-1})(3)$.
  \item Υπολογίστε το $\dint_0^4 f^{-1}(y)\,dy$, αν γνωρίζετε ότι $\dint_0^5 f(x)\,\d x = 9$.
  \item Αν η $f$ είναι κοίλη στο $[0, 5]$, τι κυρτότητα έχει η $f^{-1}$; Αιτιολογήστε γραφικά και αλγεβρικά.
  \item Ένα σημείο κινείται στην $C_{f^{-1}}$ με τη τεταγμένη να αυξάνεται με ρυθμό $3$ $\mu/\sec$. Τη στιγμή που η τεταγμένη είναι $y_0 = 3$, να βρείτε πόσο γρήγορα μεταβάλλεται η τετμημένη.
\end{erwthma}
\end{thema}
